\section{The Enum Declared Twice}
\label{sec:ch13-the-enum-declared-twice}
\index{enum!declared in two subgraphs|(}%

The \code{ProductCategory} enum has five members and is declared in two
services. Catalog owns \code{Product.category}; Pricing declares an
\code{@external} copy of that field because \code{shippingCost} requires it, and
declaring the field means declaring the enum. Nobody chose to duplicate it. It
arrived with the directive.

\index{composition!checks structure}%
Chapter~\ref{ch:composition} already has a case where the two copies disagree and
the composer refuses, so the obvious question is answered. The interesting
question is what \enquote{disagree} means, and the answer is that it depends on
something neither declaration mentions.

\subsection*{Both directions at once}

Add a member to Pricing's copy and the composer produces the longest single
message in this book:

\begin{minted}[fontsize=\scriptsize,breaklines]{text}
Enum "ProductCategory" was used as both an input and output but was inconsistently defined across inclusive subgraphs.
To update an Enum used as both an input and output, add any new Enum values with the @inaccessible directive in the
origin subgraph. Next, add those new Enum values to all other subgraphs that define the Enum—this time without the
@inaccessible directive. Finally, once all subgraphs have been updated, remove @inaccessible from the Enum values in
the origin subgraph.
\end{minted}

Remove a member instead and the message is identical, which is a small point
against it and a large point in its favour: it is not describing your edit, it is
describing the constraint.

The constraint is in the first line. The enum is the return type of
\code{Product.category}, and it is also an argument type, inside the filter
input that chapter~\ref{ch:data-without-the-n-plus-1} generated when it added
filtering:

\begin{graphqlsdl}[fontsize=\footnotesize]
input ProductCategoryOperationFilterInput {
  eq: ProductCategory
  neq: ProductCategory
  in: [ProductCategory!]
  nin: [ProductCategory!]
}
\end{graphqlsdl}

A
value travelling out and a value travelling in want opposite things from a merge,
so there is no safe answer and the composer declines to pick one. Three of those
five lines are then a migration recipe, which I have not seen a composer offer
anywhere else and which is genuinely useful: mark the new value
\code{@inaccessible} where it is added, roll it out everywhere, then take the
directive off.

\subsection*{One direction, and the two different answers}

\index{enum!output position, values unioned}%
Now the half that is easy to miss. Take a new enum that is only ever returned.
Two services declare it with different members:

\begin{graphqlsdl}[fontsize=\footnotesize]
enum StockCondition { NEW  REFURBISHED }   # inventory
enum StockCondition { NEW  DAMAGED }       # pricing
\end{graphqlsdl}

This composes, silently, and the client-facing enum is the union:

\begin{graphqlsdl}[fontsize=\footnotesize]
enum StockCondition { NEW  DAMAGED  REFURBISHED }
\end{graphqlsdl}

\index{enum!input position, values intersected}%
Now the same two declarations used only as an argument. The pair has to change,
because three of Mosaic's seven subgraphs have no root field to hang an argument
on and Pricing is one of them, so the committed case declares the enum in
Inventory and Accounts instead. The members are the same two sets:

\begin{graphqlsdl}[fontsize=\footnotesize]
enum StockCondition { NEW  REFURBISHED }   # inventory
enum StockCondition { NEW  DAMAGED }       # accounts
\end{graphqlsdl}

That composes silently too, and the client-facing enum is this:

\begin{graphqlsdl}[fontsize=\footnotesize]
enum StockCondition { NEW }
\end{graphqlsdl}

Both merges are correct and I would defend either. Every value some service can
emit has to be nameable by a client reading, so an output enum is unioned. Only
values every service accepts can be safely offered to a client writing, so an
input enum is intersected. What is uncomfortable is the second one's failure
mode. Inventory has a field that takes \code{REFURBISHED}, that field is in the
public graph, and \code{REFURBISHED} is not, because a service in another domain
happens to declare an enum of the same name without it. Nothing was reported at
any severity. Somebody's feature left the API because somebody else's schema was
merged with theirs.

\index{enum!which merge you get depends on use}%
Which leaves a rule that is easy to state and hard to remember: \textbf{the
composer's treatment of an enum is decided by where the enum is used, not by how
it is declared}. The same two files compose three different ways depending on
whether the type appears in a return position, an argument position, or both. And
the position is a property of the whole graph rather than of either declaration,
so it can change because of a schema you have never read. Add a filter argument
to your enum in one service and an unrelated service's spare member becomes a
composition error.

\medskip
So enums are checked, and checked in a way that repays understanding. The
question the last section asks is what happens when there is nothing left to
check.

\index{enum!declared in two subgraphs|)}%
